% Figure: temperature-entropy (T-S) diagram of a Carnot cycle. Two isotherms
% (horizontal) and two adiabats (vertical) bound a rectangle. The enclosed
% area is the net work; the area under the top isotherm is the heat in.
\begin{figure}[htbp]
\centering
\begin{tikzpicture}
\begin{axis}[
  width=11cm, height=8cm,
  xlabel={entropy $S$ (J/K)}, ylabel={temperature $T$ (K)},
  xmin=0.5, xmax=4.5, ymin=200, ymax=650,
  axis lines=left,
  tick label style={font=\scriptsize},
  label style={font=\small},
  clip=false,
]
% Carnot rectangle: hot isotherm at T_H=600 from S=1 to S=4; cold isotherm
% at T_C=300 from S=4 back to S=1; adiabats are the vertical sides.
\addplot[engblue, very thick] coordinates {(1,600) (4,600)};
\addplot[engblue, very thick] coordinates {(4,600) (4,300)};
\addplot[engblue, very thick] coordinates {(4,300) (1,300)};
\addplot[engblue, very thick] coordinates {(1,300) (1,600)};
% shade enclosed work area
\draw[englime!30!white, opacity=0.45, fill] (axis cs:1,300) rectangle (axis cs:4,600);
% redraw boundary on top of fill
\addplot[engblue, very thick] coordinates {(1,600) (4,600) (4,300) (1,300) (1,600)};
% labels
\node[engblue, font=\scriptsize, anchor=south] at (axis cs:2.5,600) {hot isotherm $T_H$, heat in $Q_H=T_H\,\Delta S$};
\node[engblue, font=\scriptsize, anchor=north] at (axis cs:2.5,300) {cold isotherm $T_C$, heat out $Q_C=T_C\,\Delta S$};
\node[enggray, font=\scriptsize, anchor=east] at (axis cs:1,450) {adiabat};
\node[enggray, font=\scriptsize, anchor=west] at (axis cs:4,450) {adiabat};
\node[font=\scriptsize, align=center] at (axis cs:2.5,450) {net work\\ $W=(T_H-T_C)\,\Delta S$};
\end{axis}
\end{tikzpicture}
\caption[Carnot cycle on a temperature-entropy diagram]{The Carnot cycle on a
temperature-entropy diagram is a rectangle. The two horizontal sides are
reversible isothermal heat exchanges at $T_H$ and $T_C$; the two vertical sides
are reversible adiabats, along which entropy is constant. The area under the
top edge, $T_H\,\Delta S$, is the heat absorbed; the area under the bottom edge,
$T_C\,\Delta S$, is the heat rejected; the enclosed area $(T_H-T_C)\,\Delta S$
is the net work. The thermal efficiency $1-T_C/T_H$ is the ratio of the
enclosed area to the area under the top edge, read directly off the diagram.}
\label{fig:vol04ch04:ts-diagram-carnot}
\end{figure}
